\section{Formal Framework}
\label{sec:framework}

\subsection{Schedules, costs, and interactions}

Let $G$ be a kernel in canonical form, the module after \texttt{mem2reg}. A stage $\pi$ is one pass invocation, and a
schedule $\sigma=(\pi_1,\dots,\pi_k)$ applies its stages in order to give $\sigma(G)$. We write $\pi_j\circ\pi_i$ for the
schedule that applies $\pi_i$ and then $\pi_j$.

Let $J$ be a cost of a module. The cost change of a schedule is
\begin{equation}\label{eq:delta}
  \Delta^J_\sigma(G)=J(\sigma(G))-J(G).
\end{equation}
For two stages, the interaction is the part of the pair's cost change that its two single changes do not explain,
\begin{equation}\label{eq:interaction}
\begin{split}
  I^J_{ij}(G)&=\Delta^J_{ij}(G)-\Delta^J_i(G)-\Delta^J_j(G)\\
  &=\Delta^J(\pi_j\mid\pi_i(G))-\Delta^J(\pi_j\mid G),
\end{split}
\end{equation}
where $\Delta^J(\pi_j\mid\pi_i(G))=J(\pi_j(\pi_i(G)))-J(\pi_i(G))$ is the change that $\pi_j$ causes after $\pi_i$. The
interaction is zero when the two passes act independently on $G$ for the cost $J$. A positive value means that the pair's
combined change is larger than the sum of its single changes.

The study uses two costs. The instruction count $N$ of the optimized module is an exact integer, measured before lowering, so
the backend does not affect it. The wall-clock time $J$ is measured with noise (Section~\ref{sec:timing}).

\subsection{Instruction-count interaction}

We normalize the instruction-count interaction by the canonical count $N(G)$,
\begin{equation}\label{eq:instcount}
  \tilde I_{ij}(G)=\frac{I^N_{ij}(G)}{N(G)},
\end{equation}
and average it over the $K$ kernels to give the tensor $\mathbf{M}^N_{\mathrm{pass}}$. Its entries are fractions of the
canonical instruction count.

\begin{remark}[Self-interaction identity]
\label{rem:identity}
Suppose that a second application of $\pi_i$ leaves the module unchanged. Then $\Delta^N_{ii}(G)=\Delta^N_i(G)$, and the
self-interaction is $\tilde I_{ii}(G)=(\Delta^N_{ii}(G)-2\Delta^N_i(G))/N(G)=-\Delta^N_i(G)/N(G)$. The identity is exact,
because the counts are integers of the same module. Section~\ref{sec:results} checks it on every kernel where the premise
holds.
\end{remark}

\subsection{Timing interaction}

The timing interaction is divided by the canonical time of the sweep's window,
\begin{equation}\label{eq:timing}
  \tilde I^t_{ij}(G)=\frac{I^t_{ij}(G)}{J_0(G)},
\end{equation}
where $I^t_{ij}$ is the interaction of the wall-clock cost. Both deltas of a pair are measured in one exclusive window, and the
noise is estimated from the sweeps whose optimized IR equals the canonical IR (Section~\ref{sec:timing}).

\subsection{The single-effect ratio baseline}

This baseline defines the transition tensor with a floor in the denominator,
\begin{equation}\label{eq:ratio}
  [\mathbf{M}^{\mathrm{ratio}}_{\mathrm{pass}}]_{ij}=\frac{1}{K}\sum_{G}\frac{I^t_{ij}(G)}{\max\bigl(\varepsilon,\;\lvert\Delta^t_j(G)\rvert+\lvert\Delta^t_i(G)\rvert\bigr)},
\end{equation}
with $\varepsilon=\MpassEpsilonNs{}\,\mathrm{ns}$. Its denominator is built from two noisy single-pass measurements. It is
unstable when both passes change the cost by about the noise level, and the floor only prevents division by zero. The paper
reports this estimator as a diagnostic and does not use it for its conclusions.

\subsection{Commutators}

The commutator of two stages on one kernel is
\begin{equation}\label{eq:commutator}
\begin{split}
  [\pi_i,\pi_j]_J(G)&=\frac{J((\pi_j\circ\pi_i)(G))-J((\pi_i\circ\pi_j)(G))}{J(G)}\\
  &=\frac{\Delta^t_{ij}(G)-\Delta^t_{ji}(G)}{J(G)}.
\end{split}
\end{equation}
It is zero when the two orders cannot be told apart at the resolution of the measurement.

\subsection{Redundant executions and compile-time waste}

The pass manager runs each pass as a separate execution. An execution is redundant when the IR printed after it is unchanged
from the IR before it. For a schedule $\sigma$, let $E(\sigma)$ be the number of pass executions and $R(\sigma)$ the number of
redundant ones. The redundant-pass ratio is
\begin{equation}\label{eq:rpb}
  \rho_{\mathrm{RPB}}(\sigma)=\frac{R(\sigma)}{E(\sigma)}.
\end{equation}
Let $P(\sigma)$ be the total pass time of the sweep's optimizer run, as its \texttt{-time-passes} report gives it. The pass
time that redundant executions consume is estimated as
\begin{equation}\label{eq:waste}
  \mathcal{W}(\sigma)=P(\sigma)\,\rho_{\mathrm{RPB}}(\sigma).
\end{equation}
This estimate assumes that every execution costs the same share of $P(\sigma)$. The telemetry does not record the time of each
execution, so the estimate is not a measurement of those times. The wall time of the optimizer process is not used, because it
also includes process start-up, parsing, and trace output.

\subsection{Objectives and the hypervolume}

Each sweep is described by three objectives, all to be maximized. Let $J_0$ be the canonical cost measured in the sweep's own
window, $T_0$ the canonical text size, and $C_0$ the compile time of the kernel's $-\mathrm{O}0$ sweep. For a sweep with cost
$J$, text size $T$, and compile time $C$, the objectives are
\begin{equation}\label{eq:objectives}
  s=\frac{J_0}{J_0+J},\qquad b=\frac{T_0}{T_0+T},\qquad c=\frac{C_0}{C_0+C},
\end{equation}
where $C$ is the optimizer wall time recorded for the sweep. Each objective is $\tfrac12$ when the sweep matches its reference,
so the reference point of the hypervolume is $\mathbf{r}=(\tfrac12,\tfrac12,\tfrac12)$. A schedule contributes only when it is
better than the reference on all three objectives. The hypervolume of a finite set $P$ is the Lebesgue measure of
$\bigcup_{p\in P}[\mathbf{r},p]$.

\subsection{What an UNSAT answer establishes}
\label{sec:unsat}

The checker encodes a fragment $\mathcal{F}$ of LLVM integer code: the operations \texttt{add}, \texttt{sub}, \texttt{mul},
\texttt{and}, \texttt{or}, \texttt{xor}, \texttt{shl}, \texttt{lshr}, \texttt{ashr}, \texttt{icmp}, \texttt{select},
\texttt{zext}, \texttt{sext}, \texttt{trunc}, and \texttt{ret}, on values of at most $64$ bits. Each instruction in
$\mathcal{F}$ may carry only the flags of Table~\ref{tab:flags}. A block containing any other instruction, operand, or
flag is not encodable, and the checker reports it as such. Each defined value $v$ is encoded as a pair $(\hat v,\bar v)$
of a bit-vector and a poison flag. A live-in is a value used but not defined in the block.

\begin{table}[t]
\centering
  \caption{Poison conditions of the modeled flags, on defined operands. Each condition is an exact predicate on the
  bit-vectors, and any instruction in the fragment whose condition holds yields poison.}
  \label{tab:flags}
  \begin{tabular}{|l|l|p{0.52\linewidth}|}
    \hline
    Instruction & Flag & Poison when \\
    \hline
    \texttt{add}, \texttt{sub}, \texttt{mul} & \texttt{nsw} & the exact result lies outside the signed range of the width \\
    \hline
    \texttt{add}, \texttt{sub}, \texttt{mul} & \texttt{nuw} & the exact result lies outside the unsigned range of the width \\
    \hline
    \texttt{shl} & \texttt{nsw} & the signed value times $2^{\text{amount}}$ lies outside the signed range \\
    \hline
    \texttt{shl} & \texttt{nuw} & the unsigned value times $2^{\text{amount}}$ lies outside the unsigned range \\
    \hline
    \texttt{lshr}, \texttt{ashr} & \texttt{exact} & any bit shifted out is one \\
    \hline
    \texttt{or} & \texttt{disjoint} & the two operands share a set bit \\
    \hline
    \texttt{zext} & \texttt{nneg} & the operand is negative as a signed value \\
    \hline
    \texttt{trunc} & \texttt{nsw}, \texttt{nuw} & the truncated value, extended back, differs from the operand \\
    \hline
    \texttt{icmp} & \texttt{samesign} & the operands have different sign bits \\
    \hline
    \hline
    \texttt{shl}, \texttt{lshr}, \texttt{ashr} & none & the shift amount is at least the bit width \\
    \hline
  \end{tabular}
\end{table}

For a name $x$ defined in both blocks, the refinement query is
\begin{equation}
  V_x=\lnot\bar{x}_s\;\wedge\;\bigl(\bar{x}_t\;\vee\;\hat{x}_s\neq\hat{x}_t\bigr),
\end{equation}
where the subscripts $s$ and $t$ denote the source and target blocks. A violation is an assignment of non-poison live-in
values for which the source defines $x$ without poison and the target does not define it equally.

\begin{lemma}[Encoding agreement]
\label{lem:agreement}
For each instruction in $\mathcal{F}$, if every operand's encoded pair agrees with the LLVM value and poison of that operand,
then the instruction's encoded pair agrees with the LLVM result. The value term applies the operation to the bit-vectors. The
poison term is the disjunction of the operands' poison bits, the shift-amount condition, and the flag conditions of
Table~\ref{tab:flags}. For \texttt{select}, the poison term depends only on the chosen arm and the condition.
\end{lemma}

The lemma has two supports. The first is by construction. Each flag condition is an exact predicate on the bit-vectors.
The overflow predicates widen the operands by one bit for \texttt{add} and \texttt{sub}, and by the full width for
\texttt{mul}. The second is by exhaustive evaluation. The production encoder is run on every concrete case for widths up to
\EncodingMaxWidth{}, with every defined value and one poisoned operand, and each result is compared with an independent
reference written from the poison rules. The check covers \EncodingCases{} cases in \EncodingFamilies{} instruction
families, with \EncodingMismatches{} disagreements (Table~\ref{tab:encoding}). The exhaustive check covers small widths only,
so the lemma for larger widths rests on the construction argument.

\begin{table}[t]
\centering
\caption{Exhaustive agreement between the Z3 encoder and the reference semantics, widths 1 to 4, per instruction family.}
\label{tab:encoding}
\setlength{\tabcolsep}{3pt}
\begin{tabular}{|l|r|r|}
\hline
Family & Cases & Mismatches \\
\hline
add & 404 & 0 \\
\hline
add nsw & 404 & 0 \\
\hline
add nsw nuw & 404 & 0 \\
\hline
add nuw & 404 & 0 \\
\hline
and & 404 & 0 \\
\hline
ashr & 404 & 0 \\
\hline
ashr exact & 404 & 0 \\
\hline
icmp & 4,040 & 0 \\
\hline
icmp samesign & 4,040 & 0 \\
\hline
lshr & 404 & 0 \\
\hline
lshr exact & 404 & 0 \\
\hline
mul & 404 & 0 \\
\hline
mul nsw & 404 & 0 \\
\hline
mul nsw nuw & 404 & 0 \\
\hline
mul nuw & 404 & 0 \\
\hline
or & 404 & 0 \\
\hline
or disjoint & 404 & 0 \\
\hline
select & 1,212 & 0 \\
\hline
sext & 28 & 0 \\
\hline
shl & 404 & 0 \\
\hline
shl nsw & 404 & 0 \\
\hline
shl nsw nuw & 404 & 0 \\
\hline
shl nuw & 404 & 0 \\
\hline
sub & 404 & 0 \\
\hline
sub nsw & 404 & 0 \\
\hline
sub nsw nuw & 404 & 0 \\
\hline
sub nuw & 404 & 0 \\
\hline
trunc & 74 & 0 \\
\hline
trunc nsw & 74 & 0 \\
\hline
trunc nsw nuw & 74 & 0 \\
\hline
trunc nuw & 74 & 0 \\
\hline
xor & 404 & 0 \\
\hline
zext & 28 & 0 \\
\hline
zext nneg & 28 & 0 \\
\hline
\end{tabular}
\end{table}


\textbf{Soundness of the query.} Let $B_s,B_t\in\mathcal{F}$ have the same live-in names and widths, and let $x$ be a name
defined in both blocks. Suppose that no assignment of non-poison values to the live-ins satisfies $V_x$. Take an assignment
under which $B_s$ defines $x$ without poison. By Lemma~\ref{lem:agreement}, applied to the instructions of each block in
definition order, the encoded pairs of every name agree with LLVM under that assignment. So $\bar{x}_s$ is false, and
$\hat{x}_s$ is the value of $x$ in $B_s$. If $B_t$ did not define $x$ without poison, or computed a different value, then
$\bar{x}_t\vee\hat{x}_s\neq\hat{x}_t$ would hold, and the assignment would satisfy $V_x$. That contradicts the assumption.
So $B_t$ defines $x$ without poison and computes the same value. This argument is the whole soundness claim, and it holds
only as far as the lemma does.

The argument makes no claim about values defined in only one block, about live-ins that differ, about poison live-ins, or
about facts that hold only in the surrounding function. The last limit matters in practice. An optimizer may add a
\texttt{nuw} flag because a range fact from another block rules out wrap-around. A block-local check does not see that
fact. It then reports a counterexample that the whole function never reaches. We therefore report a SAT answer as a
candidate counterexample, and we separate it from behavior observed by execution.

\subsection{Exact hypervolume}

\textbf{Exactness of the algorithm.} The hypervolume is computed with exact rational arithmetic over the stored binary64
objective values. The routine \texttt{hypervolume\_\allowbreak{}3d} first discards the points that are not strictly greater than the
reference point in every coordinate. A discarded point has a box of zero measure, so discarding it does not change the
union. The three-dimensional union is then cut into slabs along the third coordinate, and the cross-section of each slab is
measured with a two-dimensional routine. That routine sorts the points by their first coordinate and adds one strip for each
point that raises the running maximum of the second coordinate. The volume is the sum of the cross-sections times the slab
thicknesses. The routine is covered by unit tests that accompany the artifact. The result is exact for the
stored objective values, not for the measurements behind them, because each objective value is a ratio of measured times.
